\documentclass[11pt,letterpaper]{article}
\usepackage[T1]{fontenc}
\usepackage[utf8]{inputenc}
\usepackage[scaled=0.95]{helvet}
\renewcommand{\familydefault}{\sfdefault}
\usepackage[top=1.7cm,bottom=1.5cm,left=2.2cm,right=2.2cm]{geometry}
\usepackage{microtype}
\usepackage[hidelinks]{hyperref}
\pagestyle{empty}
\setlength{\parindent}{0pt}
\setlength{\parskip}{7pt}

\begin{document}

{\large\bfseries Grover Aruquipa}\par\vspace{-6pt}
{\small Qu\'ebec, QC, Canada \ \textbullet\ \ grover.aruquipa.9@gmail.com \ \textbullet\ \ groveraruquipa.github.io}

\vspace{4pt}
Hiring Team, Amazon Robotics\\
Mechanical Engineer II

Dear Hiring Team,

My first robot, during my bachelor's in Bolivia, was a reconfigurable Delta robot for product picking. Since then almost everything I have built has been about the same question: how to pick an object quickly and reliably with a mechanism that is simple to build. Amazon Robotics runs that question at a scale no lab can, and I would like to work on it with you.

I am finishing my PhD at Universit\'e Laval with Prof.\ Cl\'ement Gosselin and Prof.\ Patrice Lambert, and I defend in early 2027. I owned three robots end to end, with 4, 7 and 8 actuators and integrated grippers. For each one I wrote the specification, chose the concept, did the SolidWorks design and the FEA, selected the motors and transmissions, built the prototype, planned and ran the validation tests, and redesigned what did not meet the target. I am now building a Delta robot with an integrated Sarrus gripper that picks fragile objects such as eggs at speed by estimating the grasp force from the motor currents, without a force sensor. My 4-DoF coaxial wrist-gripper turned its tool ten full revolutions with less than one degree of tracking error while holding the object.

I am comfortable when the requirements change, because research hardware is mostly that. I design for the process that will make the part, I build proof-of-concept models fast, and I let the test data decide. I also work across disciplines: I wrote the 500\,Hz real-time control of my robots myself, and before the PhD I worked in industry as a computer vision engineer at Sencon and as a data science intern at Hitachi ABB Power Grids in Zurich. When a system fails, I can usually tell whether the cause is mechanical, electrical or in the software.

I mentor two graduate students day to day on an 8-actuator robot and I am first author of three journal papers with four more under review, so I am used to owning a project, explaining it clearly and delivering it on a deadline.

I am available from January 2027 and open to relocation and to travel as the program requires. Prof.\ Gosselin, Prof.\ Lambert and Univ.-Prof.\ Andreas M\"uller of JKU Linz can speak to my work; their details are in my CV.

Thank you for your time. I would be glad to walk you through my robots on video or in person.

Sincerely,\\
Grover Aruquipa

\end{document}
